\clearpage
\hypertarget{gl-data}{}
\section*{链接注释 · 数据与模型}
{\small 词条后的数字通往来源，返回链接指向对应正文。另读 \hyperlink{gl-eval}{统计与指标}、\hyperlink{gl-method}{研究设计与输出}。}

\gloss{retrieval}{卫星降水反演}{用卫星测到的辐射等信息推断对应时空范围的降水。未来时段预报还要定义预报提前量与未来目标。}{\src{P2}\src{R3}\back{question}{第1节}}
\gloss{latency}{近实时与时延}{从观测结束到可用产品输出的耗时应可核验，涉及数据到达、处理、缺报和发布。历史有效时间匹配不能单独证明业务时效。}{\src{P2}\back{question}{第1节}}
\gloss{imerg}{IMERG V07 Final}{NASA 的 GPM 多卫星融合降水产品。V07 是版本，Final 是事后处理流，带月尺度雨量站资料调整。参考产品存在自身误差。}{\src{R1}\src{P1}\back{data}{第2节}}
\gloss{time}{时间支撑与 UTC}{时间支撑表示测量实际覆盖的区间。UTC 是统一时间基准。目标窗左右端是否纳入、观测有效时间与文件获取时间都应记录。}{\src{P2}\back{data}{第2节}}
\gloss{satellite}{光谱通道、亮温与时槽}{通道对应不同光谱测量，亮温是将辐射换算成等效黑体温度。时槽是名义采样位置，真正的观测起止仍需解码核对。云顶信号与地面雨强的关系依赖天气环境。}{\src{R3}\src{P2}\back{data}{第2节}}
\gloss{gfs}{GFS、f006 与资料版本时效}{GFS 是全球数值天气模式。f006 表示距起报 6 小时，预报有效时刻通常为起报加时效。可用时间还要核查生产与传输，因此需要实际可用档案版本。}{\src{R2}\src{P3}\back{data}{第2节}}
\gloss{grid}{格点与 SP04}{SP04 为项目的空间绑定协议。复算需匹配坐标参考、经纬轴、像元位置及监督范围。经纬度网格的同度数面积随纬度变。}{\src{P1}\src{P2}\back{data}{第2节}}
\gloss{mask}{掩膜}{布尔资格标记。云南监督掩膜定空间范围，数据质量掩膜排除无效值，地下层掩膜排除对应地表以下的气压层。比较使用共同有效格点。}{\src{P3}\back{data}{第2节}}
\gloss{scaler}{标准化}{常见写法 \(x'=(x-\mu)/\sigma\)。统计量用训练数据拟合，按同一字段、单位和身份应用到后续划分，避免开发或测试信息进入训练。}{\src{P2}\src{P3}\back{data}{第2节}}
\gloss{dem}{DEM 与地形梯度}{DEM 是海拔栅格。以水平米制距离求 \(\nabla h=(\partial h/\partial x,\partial h/\partial y)\)，单位为 m/m。应记录投影、聚合、缺失与接缝。}{\src{P3}\back{data}{第2节}}
\gloss{dote}{动态地形强迫与 DOTE}{可用 \(u\,\partial h/\partial x+v\,\partial h/\partial y\) 理解沿地形的运动表征，风用 m/s 时结果单位为 m/s。正值表示沿所定义坐标吹向上坡。项目特征还绑定物理场与掩膜，具体实现见工程记录。}{\src{P3}\src{P6}\back{model}{第3节}}
\gloss{mee}{MEE 与物理条件化}{项目环境分支把 GFS 物理字段作为条件输入，并验证变量、气压层、单位和掩膜。显式守恒约束应说明守恒方程、离散方式与检验。}{\src{P3}\src{P4}\back{model}{第3节}}
\gloss{unet}{U-Net 与跳接}{编码阶段聚合上下文，解码阶段恢复空间表达。跳接把相应尺度的空间信息传至解码层。项目借用其结构思想做格点估计。}{\src{R4}\back{model}{第3节}}
\gloss{heads}{发生头与条件分位头}{共用特征后分别输出有雨概率与有雨条件下的分位节点。条件头在参考有雨样本上直接接受分位监督。}{\src{P2}\back{model}{第3节}}

\clearpage
\hypertarget{gl-eval}{}
\section*{链接注释 · 统计与指标}
\gloss{quantile}{分位数与条件分位数}{\(Q_\tau=\inf\{r:F(r)\ge\tau\}\)。条件分位数使用指定条件下的分布。本项目 \(\tau_i=(i-0.5)/32\)，最高为 0.984375。例：若模型校准良好，按同样条件汇总时，约 90\% 的观测应不超过预测的 0.9 分位数，离散值的并列需要注明。}{\src{P2}\src{R5}\back{model}{第3节}}
\gloss{pinball}{Pinball 损失}{令 \(e=y-q_\tau\)，则 \(\rho_\tau(e)=e(\tau-\mathbf1[e<0])\)。在 \(\tau=0.9\) 时，低估 2 单位得 1.8，高估 2 单位得 0.2。案例同时报告 \(\log(1+R)\) 域和物理域。物理 Pinball/20 按引用报告的固定尺度归一化。报告样本、分位水平、尺度与分母。}{\src{P2}\src{P7}\src{R5}\back{diagnosis}{第5节}}
\gloss{brier}{Brier 分数}{对样本计算 \((p_i-z_i)^2\) 后取平均，\(z_i=1\) 表示事件发生。分数越低越好。例如预测 0.8，若事件发生得 0.04；若没发生得 0.64。}{\src{R6}\back{diagnosis}{第5节}}
\gloss{rank}{AP、AUROC 与正例定义}{AP 按召回增加加权汇总精确率。AUROC 汇总阈值变化时真阳性率与假阳性率的关系，常可解释为正例得分超过负例的概率，并列按半计。强雨稀少时，报告正例比例和 AP，并结合固定阈值计数。换正例定义或样本会改变可比性。}{\src{P8}\src{R7}\back{diagnosis}{第5节}}
\gloss{confusion}{命中、误报、漏报和阈值指标}{TP 为真强雨且报强雨，FP 为非强雨却报强雨，FN 为真强雨却没报。\(POD=TP/(TP+FN)\)，精确率 \(=TP/(TP+FP)\)，误报比 \(=FP/(TP+FP)\)，假阳性率 \(=FP/(FP+TN)\)。分母为零时明确报告处理约定。}{\src{P8}\back{diagnosis}{第5节}}
\gloss{tail}{覆盖、节点越线与尾部识别}{覆盖比较真值 \(R\) 与节点 \(q_\tau\)，例如计 \(R\le q_{32}\)；节点越线比较 \(q_{32}>20\)。固定点值检出则比较读出 A 与 20。有限节点的兼容概率界限反映信息不足，抽样置信区间反映估计不确定性。}{\src{P6}\src{P7}\back{diagnosis}{第5节}}
\gloss{readout}{读出、A 点值与 \(P20\)}{读出把模型输出转成可比较的量。固定 A 为发生概率乘以32 个物理分位节点的算术平均。若有雨定义 \(R>0.1\)，精确的全分布均值还会含低雨强条件项。\(P20\) 指超过 20 的概率，与 A 的物理雨强值有不同单位和含义。}{\src{P2}\src{P7}\back{model}{第3节}}
\gloss{calibration}{概率校准}{在概率近似 0.2 的一组独立可比样本中，实际发生率应接近 0.2。排序相同也可有不同校准。概率可靠性需要足量样本、置信信息及清晰的拟合与评价划分。}{\src{R6}\src{P7}\back{diagnosis}{第5节}}
\gloss{block}{配对与区块重采样}{先在同一单位上做候选减对照，再按事前确定的事件或时间块取样以保留块内相关性。区块过少、块间强依赖或覆盖偏倚时，明确推断受限。}{\src{P1}\back{experiment}{第4节}}
\gloss{event}{像元时、场景与天气过程}{像元时是格点乘目标窗，场景是一整个匹配格点场。过程是满足事前时空连通规则的一组降水暴露。应区分完整、匹配和可评价的过程。}{\src{P1}\src{P6}\back{verification}{第6节}}

\clearpage
\hypertarget{gl-method}{}
\section*{链接注释 · 研究设计与输出}
\gloss{split}{Train、Dev、Final Test 与泄漏}{Train 拟合参数，Dev 支持选择与诊断，Final Test 用于冻结方案后的评估。使用最终测试指标来挑候选、阈值或数据处理规则，会把测试信息泄露给开发。时间留出检验仍受参考来源限制。}{\src{P1}\src{P2}\back{question}{第1节}}
\gloss{seed}{随机种子}{控制初始化、抽样和顺序等随机过程。相同 seed 有助于配对，多个 seed 描述训练随机性。恢复一次运行不增加独立种子，多个随机种子也不增加独立天气过程。}{\src{P4}\back{question}{第1节}}
\gloss{ablation}{消融与匹配对照}{有意去掉或替换某种输入、模块或损失，用相同评价协议比较。若同步改变容量、数据或训练时长，应明确结果混入了这些因素。真实/归零对照是输入消融的一种。}{\src{P4}\src{P6}\back{model}{第3节}}
\gloss{population}{样本集合与共同有效样本}{比较所用样本必须有可核验的时空身份。相同场景数并不自动表示真值、格点、时标和掩膜相同。先核对集合，再比较指标。}{\src{P1}\src{P8}\back{experiment}{第4节}}
\gloss{guardrail}{主指标、保护线与停止规则}{主指标回应研究问题。保护线限定可接受代价，如弱雨误差或误报。停止规则规定候选何时封存或回到诊断，阈值应在见到结果前写定。}{\src{P7}\src{P8}\back{experiment}{第4节}}
\gloss{checkpoint}{检查点、E9 与指纹}{检查点保存某一训练状态。E9 是第九轮的固定终点，和事后挑最佳轮次的规则不同。文件指纹用于确认身份，还需数据与评估记录才能回溯结果。}{\src{P4}\src{P5}\back{experiment}{第4节}}
\gloss{reference}{参考一致性与独立验证}{同源参考评价沿用与训练标签相同的产品来源。独立验证要求检查待验证模型与观测的来源链、共享订正与覆盖，并在可比时空支撑上配对。}{\src{P1}\src{R1}\back{verification}{第6节}}
\gloss{qpe}{雷达 QPE 与站点代表性}{QPE 是定量降水估计。雷达反射率到地面雨强可能受遮挡和订正影响，站点的点测与格点面积平均存在尺度差异。}{\src{P9}\back{verification}{第6节}}
\gloss{dpr}{GPM DPR}{双频降水雷达。过境足迹能提供局部观测，瞬时空间足迹与半小时面积平均须按评价协议处理，并核对资料源链。}{\src{P9}\src{R1}\back{verification}{第6节}}
\gloss{claim}{主张、结果与科学验收}{结果是某口径下的观测差值，主张增加了范围和解释。工程完成、开发筛查、机制归因与泛化各有证据条件。正文应同时给出关键反例。}{\src{P1}\src{P4}\back{communication}{第7节}}

\clearpage
\hypertarget{sources}{}
\section*{证据与延伸阅读}
\newcommand{\repo}{https://github.com/flyeagle-cex/YunTAPR-Net/blob/23900b20579afb408727f9a7f2aaf0b4d286f712}
{\small 讲义从 Obsidian 的项目主页、研究设计与证据地图、四项创新、山区强雨、TAIL20、A20、IMERG 参考评价快照及文献索引组织而成。下列项目链接指向可复核的固定提交和公开资料。}

\sub{项目记录}
\source{P1}{\repo/docs/research_readiness/20261121_stage45_observation_feasibility/IMERG_REFERENCE_EVENT_PROTOCOL_V1.md}{IMERG 参考评价与事件协议 V1}
\source{P2}{\repo/docs/phase_b_v2_protocol_candidates/v1_20261010/METHODS_AND_EXPERIMENTAL_DESIGN.md}{数据时间绑定、训练划分和概率输出协议}
\source{P3}{\repo/docs/research_readiness/20261115_core_innovation_repair/OBSIDIAN_SNAPSHOT.md}{四项创新主线与物理输入修复快照}
\source{P4}{\repo/docs/research_readiness/20261025/B0_B4_STAGE_ACCEPTANCE.md}{B0--B4 阶段交接与匹配比较}
\source{P5}{\repo/docs/research_readiness/20261110_cosine9/PROTOCOL.md}{COSINE9 固定九轮比较协议}
\source{P6}{\repo/docs/research_readiness/20261115_mountain_pair_results/OBSIDIAN_TERMINAL_SNAPSHOT.md}{REAL/ZERO 山区强雨配对结果}
\source{P7}{\repo/docs/research_readiness/20261119_conditional_intensity_diagnostic/STRONG_RAIN_DETECTION_EVIDENCE.md}{TAIL20 与强雨检出瓶颈}
\source{P8}{\repo/docs/research_readiness/20261119_conditional_intensity_diagnostic/A20_POINT_RANK_RESULTS.md}{A20 与 COSINE9 配对排序结果}
\source{P9}{\repo/docs/research_readiness/20261121_stage45_observation_feasibility/STAGE45_EVENT_MATCHING.md}{独立观测与事件配对的可行性口径}

\source{P10}{\repo/docs/research_readiness/20261121_stage45_observation_feasibility/STAGE2_EVALUATION_FREEZE_V1.md}{冻结的开发评价、指标与最终测试边界}
\source{P11}{\repo/STAGE_ACCEPTANCE.md}{地形分支与各试验阶段验收}

\source{P12}{\repo/docs/research_readiness/20261110_cosine9/SAVED_TRADEOFF.json}{COSINE9 的雨强分层代价收据}

\source{P13}{\repo/docs/research_readiness/20261119_conditional_intensity_diagnostic/A20_RESULTS.md}{A20 原固定阈值验收与九项筛查}

\sub{专业知识原始来源}
\source{R1}{https://gpm.nasa.gov/sites/default/files/2023-07/IMERG_TechnicalDocumentation_final_230713.pdf}{NASA · IMERG V07 技术文档}
\source{R2}{https://emc.ncep.noaa.gov/emc/pages/numerical_forecast_systems/gfs.php}{NOAA/NCEP · Global Forecast System}
\source{R3}{https://amt.copernicus.org/articles/15/6907/2022/}{Pfreundschuh 等（2022）· Hydronn 巴西近实时卫星降水反演}
\source{R4}{https://arxiv.org/abs/1505.04597}{Ronneberger 等（2015）· U-Net 原论文}
\source{R5}{https://people.eecs.berkeley.edu/~jordan/sail/readings/koenker-bassett.pdf}{Koenker 与 Bassett（1978）· Regression Quantiles}
\source{R6}{https://sites.stat.washington.edu/raftery/Research/PDF/Gneiting2007jasa.pdf}{Gneiting 与 Raftery（2007）· 严格适当评分规则}
\source{R7}{https://ftp.cs.wisc.edu/machine-learning/shavlik-group/davis.icml06.pdf}{Davis 与 Goadrich（2006）· 精确率—召回率与 ROC}
\medskip
\hyperlink{nav}{返回学习路径 $\rightarrow$}
